\chapter{LITERATURE REVIEW}

Efforts to automatically recognize hand-drawn circuit diagrams date back to the early 2000s. Edwards and Chandran \cite{edwards2000machine} were among the first to tackle this problem; their system identified circuit components and then worked out node segments from spatial relationships. It managed 86\% accuracy on component recognition and 92.5\% on node identification, tracing wire connections through hard-coded heuristic rules. The reliance on hand-crafted rules made the system fragile, though, and it struggled to cope with varied handwriting or unfamiliar circuit layouts. Around the same period, Ramel et al.\ \cite{ramel2000structured} tackled a related challenge in technical document recognition by splitting graphical primitives from textual annotations and matching structures against predefined symbol templates. Belongie et al.\ \cite{belongie2002shape} contributed shape context descriptors that computed histograms of relative contour point positions, producing descriptors tolerant of minor deformations, a trait especially handy for hand-drawn symbols. These early methods, for all their contributions, could not learn discriminative features from data on their own, which is what eventually drove the field toward deep learning.

Detecting and classifying individual circuit components sits at the heart of any recognition pipeline, and deep learning-based object detectors now dominate this space. The R-\gls{cnn} lineage brought region proposals into the picture: Girshick et al.\ \cite{girshick2014rich} put forward R-\gls{cnn} with selective search for candidate regions, Fast R-\gls{cnn} \cite{girshick2015fast} then shared convolutional features across those proposals, and Faster R-\gls{cnn} \cite{ren2015faster} swapped selective search out for a learned Region Proposal Network. He et al.\ \cite{he2017mask} built on this further with Mask R-\gls{cnn}, bolting on instance segmentation. Bayer et al.\ \cite{bayer2023instance} applied Mask R-\gls{cnn} to hand-drawn circuit diagrams and reached roughly 94\% mask accuracy, though they conceded that symbol recognition remained ``reasonable, yet incomplete,'' especially where handwriting was erratic. On the single-shot side, Liu et al.\ \cite{liu2016ssd} devised SSD, which tiles default bounding boxes at several scales to skip the proposal stage entirely. The \gls{yolo} family has had an outsized influence here. Redmon et al.\ \cite{redmon2016yolo} recast detection as one unified regression problem in the original \gls{yolo}. Subsequent iterations added batch normalization and anchor boxes in YOLOv2 \cite{redmon2017yolo9000}, feature pyramid predictions at multiple scales in YOLOv3 \cite{redmon2018yolov3}, Mish activation with Cross-Stage Partial connections in YOLOv4 \cite{bochkovskiy2020yolov4}, and a widely adopted PyTorch codebase in YOLOv5 \cite{jocher2020yolov5}. Wang et al.\ \cite{wang2023yolov7} then introduced YOLOv7, which brought E-\gls{elan} and model re-parameterization to the table, topping all prior real-time detectors on MS COCO in both throughput and accuracy. YOLOv7 is the architecture this project uses for component detection.

Adnan et al.\ \cite{adnan2021cnnbased} built a two-stage \gls{cnn}-based classifier in 2021 that could sort 20 distinct hand-drawn circuit components, posting 97.33\% recognition accuracy. Reddy and Panicker \cite{reddy2021handdrawn} paired YOLOv5 with probabilistic Hough transform-based node recognition; their setup hit a \gls{map} of 98.2\% for component detection and 80\% for overall schematic generation, although labeled components were left unhandled. Abhinav et al.\ \cite{abhinav2020electric} took a different route, first detecting and classifying components, then piping the resulting labeled strings into an \gls{fsm} that reconstructed the circuit.

Being able to pull component values like ``10k$\Omega$'' or ``5V'' off hand-drawn annotations matters enormously for generating netlists that actually work. Shi et al.\ \cite{shi2017crnn} put forward the \gls{crnn} architecture, which folds convolutional feature extraction together with bidirectional \gls{lstm} sequence modeling and \gls{ctc}-based transcription. That framework became the go-to for scene text recognition and underpins the \gls{ocr} module used in this project. The \gls{ctc} loss function itself was originally devised by Graves et al.\ \cite{graves2006ctc} for labeling unsegmented sequences; its chief virtue is that it sidesteps the need for pre-segmented training data, which is particularly important for handwritten text where character boundaries are blurry at best. Baek et al.\ \cite{baek2019scene} ran a thorough cross-comparison of scene text recognition models and found that attention-based decoders tend to beat \gls{ctc}-based ones on irregular text, while \gls{ctc} holds its own on regular text. Since circuit annotations are usually written horizontally, a \gls{ctc}-based \gls{crnn} was a sensible pick for this project. On the transformer front, TrOCR \cite{li2021trocr} has shown impressive results on handwritten text by capitalizing on pre-trained vision transformers, but the computational overhead it demands makes it a heavier choice for lightweight deployments like this one.

Wire detection and connectivity inference is probably the trickiest part of recognizing circuit diagrams, because wires can run in any direction, overlap one another, cross without connecting, or converge at junctions. The classical Hough transform by Duda and Hart \cite{duda1972hough} remains a staple for straight-line detection, and Matas et al.\ \cite{matas2000robust} sped it up with a probabilistic variant that samples random subsets of edge points. Reddy and Panicker \cite{reddy2021handdrawn} applied the Probabilistic Hough Transform to hand-drawn circuits but ran into trouble with curved and overlapping wires. Morphological thinning offers another avenue; Zhang and Suen \cite{zhang1984thinning} proposed a parallel thinning algorithm that collapses binary regions to single-pixel-wide skeletons without breaking topological connectivity. Hemker et al.\ \cite{hemker2024schematics} went with a three-stream pipeline, dedicating one stream to component detection, a second to line and junction detection, and combining their outputs into a \gls{json} representation that feeds netlist generation. Kelly and Cole \cite{kelly2024digitizing} managed 86.4\% digitization success and 99.6\% \gls{ocr} accuracy on printed schematics, but because their model was trained solely on computer-generated diagrams, it cannot accommodate hand-drawn variability. A different angle treats wire tracing as a graph search: the binarized image becomes a grid graph and breadth-first search or $A^{*}$ hunts for paths between component terminals. This project settled on that path-finding strategy after early Hough transform experiments fell short.

Casting a circuit as a graph, with nodes standing for components or connection points and edges for electrical links, opens the door to topological reasoning and eventual simulation. Kipf and Welling \cite{kipf2017gcn} showed how convolutions could be generalized to graph-structured data through Graph Convolutional Networks. Yamakaji et al.\ \cite{yamakaji2023circuit2graph} took this idea into the circuit domain with Circuit2Graph, converting diagrams into graph structures and hitting 97.9\% accuracy on circuit classification via \glspl{gnn}. Their work demonstrated that graph neural networks can capture the global structure of a circuit in situations where visual features alone come up short.

For most circuit recognition systems, the endgame is producing a netlist that a simulator can actually consume. Conventional \gls{eda} packages like KiCad and \gls{ltspice} let users export netlists from manually drawn schematics, yet none of them ingest a raw image and turn it into a netlist directly \cite{kicad2024spice}. Sertdemir et al.\ \cite{sertdemir2022img2sim} took a stab at closing that gap with Img2Sim. Their artificial neural network-based pipeline picks out components, works out how they are connected, and emits a full \gls{spice} netlist at better than 90\% topology accuracy. The catch is that it needs reasonably clean input images and only recognises a fixed roster of components. The lineage of \gls{spice} itself traces back to Nagel and Pederson \cite{nagel1973spice} at UC Berkeley, whose work set the convention that still anchors circuit simulation today; modern descendants like Ngspice handle transient analysis along with \gls{dc} and \gls{ac} sweeps, and commercial offerings such as PSpice and Simulink occupy the higher end of that ecosystem. For this project, CircuitJS \cite{falstad2026} was chosen as the simulation backend. Being browser-based, openly accessible, and capable of animating voltages and currents in real time made it a natural fit.

The arrival of \glspl{llm} has cracked open a new set of possibilities for explaining and reasoning about circuits. Roughly 115,000 question-image pairs spanning 5,725 distinct circuits and around 70 symbol classes are bundled in the CircuitVQA dataset \cite{li2024circuitvqa}; transformer-based \gls{vqa} models trained on it outperform generic counterparts once fine-tuning on domain-specific material kicks in. Vungarala et al.\ \cite{vungarala2024spicepilot} took a more code-centric angle with SPICEPilot, which leans on GPT-style models to write and interpret \gls{spice} code. They also flagged a recurring weakness: out-of-the-box \glspl{llm} routinely violate domain conventions unless explicitly adapted. For the conversational layer in this project, DeepSeek serves as the \gls{llm}-backed assistant that fields user questions about circuits. The pedagogical motivation comes from Chen et al.\ \cite{chen2011efficacy}, whose study of simulation-based learning in electronics education reported substantial gains in student comprehension when interactive simulations were brought into the loop. That finding shaped CIRCUITRON's framing as something more than a recognition utility, namely an educational platform.

Even though meaningful progress has accumulated across component recognition, value extraction, connectivity inference, netlist generation, and interactive Q\&A, the existing landscape stays fragmented. Most published systems handle only a slice of the pipeline. Several of them \cite{kelly2024digitizing, hemker2024schematics} target printed schematics and simply do not transfer to hand-drawn input. None of the surveyed systems output an editable, interactive schematic. \gls{vqa} and \gls{llm}-based tooling has yet to be folded into recognition pipelines in any serious way. A number of approaches \cite{reddy2021handdrawn, bayer2023instance} also leave handwritten component values entirely on the table. CIRCUITRON sets out to plug these gaps with a single end-to-end pipeline that takes a hand-drawn circuit image and gives back a simulation-ready digital schematic with \gls{llm}-based assistance baked in. The capabilities of the reviewed systems are laid out side by side in Table~\ref{tab:litreview_comparison}.

\begin{table}[htbp]
\centering
\caption{Comparison of Existing Circuit Recognition Systems}
\label{tab:litreview_comparison}
\renewcommand{\arraystretch}{1.2}
\small
\begin{tabular}{|p{2.8cm}|c|c|c|c|c|c|}
\hline
\textbf{System} & \textbf{Hand-} & \textbf{Component} & \textbf{Value} & \textbf{Wire} & \textbf{Netlist/} & \textbf{AI} \\
& \textbf{drawn} & \textbf{Detection} & \textbf{Extraction} & \textbf{Detection} & \textbf{Simulation} & \textbf{Assist} \\
\hline
Edwards \cite{edwards2000machine} & Yes & Yes & No & Yes & No & No \\
\hline
Adnan et al.\ \cite{adnan2021cnnbased} & Yes & Yes & No & No & No & No \\
\hline
Kelly \& Cole \cite{kelly2024digitizing} & No & Yes & Yes & Yes & No & No \\
\hline
Reddy \& Panicker \cite{reddy2021handdrawn} & Yes & Yes & No & Yes & Yes & No \\
\hline
Bayer et al.\ \cite{bayer2023instance} & Yes & Yes & No & Yes & No & No \\
\hline
Hemker et al.\ \cite{hemker2024schematics} & No & Yes & Yes & Yes & Yes & No \\
\hline
Img2Sim \cite{sertdemir2022img2sim} & Yes & Yes & No & Yes & Yes & No \\
\hline
Circuit2Graph \cite{yamakaji2023circuit2graph} & No & Yes & No & No & No & No \\
\hline
SPICEPilot \cite{vungarala2024spicepilot} & No & No & No & No & Yes & Yes \\
\hline
\textbf{CIRCUITRON} & \textbf{Yes} & \textbf{Yes} & \textbf{Yes} & \textbf{Yes} & \textbf{Yes} & \textbf{Yes} \\
\hline
\end{tabular}
\end{table}